

\documentclass[11pt]{article}

\usepackage[margin=1in]{geometry}
\usepackage{amsmath,amssymb,amsthm,mathtools}
\usepackage{graphicx}
\graphicspath{{figs/}}
\usepackage{booktabs}
\usepackage{tabularx}\usepackage{array}
\newcolumntype{Y}{>{\raggedright\arraybackslash}X}
\usepackage{enumitem}
\usepackage{hyperref}
\usepackage{microtype}
\newtheorem{definition}{Definition}
\newtheorem{theorem}{Theorem}
\newtheorem{lemma}{Lemma}
\newtheorem{corollary}{Corollary}
\newtheorem{proposition}{Proposition}
\newtheorem{remark}{Remark}

\newcommand{\E}{\mathbb{E}}
\newcommand{\one}{\mathbf{1}}

\title{Committee Contact-Set Privacy in Anonymous DHT Lookups:\\MUCC, Availability Floors, and Abort-Safe Equalization}
\author{}
\date{March 21, 2026 (archive v495)}


\begin{document}
\maketitle

\begin{abstract}
Anonymous DHT designs often seek to hide the destination key by making observed contact sets
``look uniform.'' We formalize this as \emph{MUCC-marginal}: the marginal probability of contacting any committee
is key-independent.
We show that MUCC implies strong and universal \emph{availability ceilings}:
to succeed with high probability, a lookup must contact a large fraction of all committees.
When threshold decoding is required and when retries must be hidden via fixed stop time (abort-safe equalization),
these ceilings become a \emph{total cover-contact floor} that cannot be avoided by splitting work into more rounds.
Our results are agnostic to the underlying anonymous routing mechanism and therefore apply as lower bounds
to a wide range of anonymous DHT and committee-based overlay designs.
The exported object is intentionally narrow: this note is the committee-contact / conditioning root for anonymous-DHT destination claims, while richer transcript equalization, monitoring manifests, notarized evaluation bundles, and release-lineage packaging remain downstream companions.
\end{abstract}

\paragraph{Compilation note.}
This \texttt{paper.tex} is intentionally concise; key deployment-mapping guidance is included as a single-page table (Section~\ref{sec:protocol-mapping}).


\paragraph{Scope.}
We isolate and formalize \emph{Committee Contact-Set Privacy in Anonymous DHT Lookups} as a narrow contract surface: a declared committee-contact witness, a MUCC-style marginal equalization claim over that witness, and the resulting availability/contact floors for successful anonymous-DHT retrieval. Within the anonymous-DHT stack, this note is the committee-contact theorem packet, not the fuller transcript-equalization, monitoring, notarization, or release note.

\paragraph{Units.} Budgets are published in bits (use $\log_2$ and $2^{\epsilon}$); results stated in nats can be converted by dividing by $\ln 2$ (see Synthesis~8).\cite{series:notation}

\paragraph{Imports.}
We treat stop-time padding and schedule compilation as imported overlay mechanisms and cite them only as black-box prerequisites.\cite{series:stoptimeaddendum,series:schedcompiler}
Trace-interface naming (committee-contact traces, exposure normal forms, ENF-ID) is imported from Synthesis~3, and repetition semantics (TW declarations, TW-ID) from Synthesis~4.\cite{series:traceifaces,series:threatwindows}
Contract objects and publication manifests follow Synthesis~0 and ABOM/OINL (Anonymity~C).\cite{series:contractobjects,series:abomoinl}
For the state-amplification viewpoint and for closed-view auditing/CPPC manifests, see the companion notes (AnonDHT~State~1 and State~3).\cite{series:state1,series:state3}
This paper does \emph{not} own full transcript-level equalization and endpoint budgeting (Anonymity~B and Anonymity~A), closed-view monitoring manifests (AnonDHT~State~3), evaluator notarization / anti-grinding evidence release (Eval~1), or public lineage / publication packaging (ABOM/OINL plus Release~A).\cite{series:profiling,series:oddsinflation,series:state3,series:evalnotary,series:abomoinl,series:release}

\paragraph{Artifact policy.}
This archive is source-only; supporting artifacts (logs, traces, scripts, certificates, and verifier outputs) are available on request.\cite{series:artifactpolicy}

\paragraph{Claim fields (non-negotiable).}
Any ``MUCC'' statement in this paper is shorthand for a claim parameterized by
(i) an exposure-normal-form id \texttt{exposure\_nf\_id} (which fixes the projection from raw lookup traces to the committee-contact witness used below) and
(ii) a threat-window id \texttt{tw\_id} (which fixes the secret, adversary model, and repetition window).
Certificates should carry both ids as claim fields; short prefixes are fine in text, but the ABOM must contain full digests.

\paragraph{Exported object and earliest sufficient stop.}
The paper exports a \emph{MUCC contact-floor packet}: the declared committee-contact witness family, the committee universe / replication / threshold / live-response parameters \((M,d,t,\alpha)\), the MUCC marginal parameter (or approximate upper-envelope slack), any fixed-stop assumptions used to hide retries, and the resulting universal availability and total-cover-contact floors.
This is the earliest sufficient packet for committee-contact claims in anonymous DHTs.
Once that packet is fixed, downstream layers may (i) strengthen it into a fuller profiling/equalization contract over richer lookup traces (Anonymity~B), (ii) monitor it behind a CPPC / closed-view auditing boundary (AnonDHT~State~3), (iii) notarize attempt logs and anti-grinding evidence for an evaluation run (Eval~1), or (iv) package the public claim surface and lineage object (ABOM/OINL plus Release~A), but those layers are out of scope here.\cite{series:profiling,series:state3,series:evalnotary,series:abomoinl,series:release}

\paragraph{Later-terse-paper citation posture.}
When a later short paper only needs to answer \emph{which committee-contact witness and MUCC contact floor were declared for this anonymous-DHT family?}, it should stop at the minimal public MUCC claim card below plus the tiny availability drill.
It should widen only when the live issue becomes richer transcript equalization (Anonymity~B), closed-view monitoring (State~3), evaluator notarization (Eval~1), or public claim / release-lineage packaging (ABOM/OINL and Release~A).\cite{series:profiling,series:state3,series:evalnotary,series:abomoinl,series:release}

\begin{table}[t]
\centering
\small
\begin{tabularx}{0.97\linewidth}{@{}p{0.30\linewidth}Y@{}}
\toprule
\textbf{Public row} & \textbf{Pinned value}\\
\midrule
Claim tuple & \texttt{exposure\_nf\_id = enf.committee\_contact.v1}; \texttt{tw\_id = tw.lookup.30d.v1} \\
Committee universe & $M=256$ committees; destination replication $d=8$; lookup succeeds after $t=4$ live destination responses \\
Live-response model & Conservative lower bound $\alpha\ge 0.8$ on a contacted destination committee returning a live response \\
Success target & Failure budget $\beta=0.05$, so the declared service goal is $\Pr[Z\ge t]\ge 0.95$ \\
MUCC witness rule & For every committee $j$ and destination set $D$, $\Pr[j\in C\mid D]=p$ (no approximate-slack relaxation claimed here) \\
Abort-safe wrapper & Three public rounds are pinned; dummy rounds continue after internal early success so the committee-contact witness is fixed-stop and replayable \\
Exported floor & The universal ceiling yields $p\ge (1-\beta)t/(\alpha d)=0.59375$, hence $k_0=Mp\ge 152$ expected contacted committees \\
Owner routing & Richer transcript equalization, CPPC monitoring, evaluator notarization, and release lineage widen into Anonymity~B, State~3, Eval~1, and ABOM/OINL + Release~A \\
\bottomrule
\end{tabularx}
\caption{A minimal public MUCC claim card. This is the smallest public answer later terse papers can quote directly when the live issue is the declared committee-contact witness and its availability / cover-contact floor rather than transcript equalization, monitoring, evaluator logs, or release packaging.}
\label{tab:minimal-mucc-card}
\end{table}

This card is intentionally non-decorative: it pins one witness projection, one threat window, one conservative liveness factor, one fixed-stop wrapper assumption, and one exported contact floor, without silently widening into profiling budgets, CPPC manifests, evaluator bundles, or release receipts.

\paragraph{One tiny availability drill.}
With the card values above,
\[
p \;\ge\; \frac{(1-\beta)t}{\alpha d}
 = \frac{0.95\cdot 4}{0.8\cdot 8}
 = 0.59375,
\]
so for $M=256$ the expected total contact floor is
\[
k_0 = Mp \ge 256\cdot 0.59375 = 152.
\]
A later terse paper may therefore say ``the committee-contact claim exported a MUCC floor of at least 152 expected committees at $M=256$'' and move on, unless the live issue is a richer transcript claim or a successor witness card.

\paragraph{One tiny fixed-stop cutover drill.}
Keep the same public card, including the three pinned public rounds and the success target $\Pr[Z\ge t]\ge 0.95$.
If an implementation tries to satisfy that card with only $48$ expected committee-contact events per round, then its total abort-safe schedule has
\[
\E[K]=3\cdot 48 = 144,
\]
which is below the declared floor $152$ from Proposition~\ref{prop:mucc-stop-time-floor}.
That design therefore fails closed under the existing MUCC card: it must either raise the total expected cover-contact to at least $152$, weaken one of the pinned service parameters $(\beta,t,\alpha,d)$, or publish a successor card.
Merely spreading an under-budget lookup across more visible rounds does not preserve the old committee-contact claim.

\paragraph{Minimal reopen ladder.}
Once the committee-contact floor itself is fixed, the next owner should be named explicitly rather than imported by implication.
If the live issue is whether the retry / dummy-fill wrapper is itself schedule-only, reopen Operational~A and the stop-time backbone rather than restating that theorem here.\cite{series:operationalA,series:stoptimeaddendum}
If the live issue is whether a deployment actually executed the pinned public schedule, reopen Operational~B for plan-pinning and epoch-conformance receipts.\cite{series:operationalB}
If the live issue is which richer transcript projection was equalized, reopen Anonymity~B; if it is public monitoring, evaluator anti-grinding evidence, or release-lineage packaging, reopen State~3, Eval~1, or Release~A respectively.\cite{series:profiling,series:state3,series:evalnotary,series:release}

\section{Introduction}
Anonymous DHT lookups aim to hide the destination key while retaining the scalability benefits of structured overlays.
A common approach is to arrange the overlay into ``committees'' or small replica sets and to make the \emph{observed contact pattern} look key-independent.
This paper isolates a particularly clean notion of this goal---\emph{marginally-uniform committee contact (MUCC) marginals}---and shows that it imposes strong and practical lower bounds.

\paragraph{Contact sets are witnesses.}
Even if payloads are protected and routes are anonymous, a partially observing adversary may learn which committees were contacted.
This ``who got touched'' footprint is a witness stream.
AnonDHT~State~1 formalizes the state-dependent amplification story~\cite{series:state1}; here we focus on the DHT-specific floor created by availability requirements.

\begin{figure}[t]
  \centering
  \fbox{\parbox{0.92\linewidth}{\centering\textbf{Figure omitted in source-only archive.}\\(Plot/diagram and scripts available on request.)}}
  \caption{Illustrative MUCC overhead curve: forcing key-independent contact marginals drives contact fraction upward as availability targets tighten.}
  \label{fig:mucc-overhead}
\end{figure}

\paragraph{Non-redundancy.}
We treat stop-time padding and schedule equalization mechanisms as imported from the series~\cite{series:stoptimeaddendum,series:schedcompiler}.
Our contribution is the \emph{availability floor} that remains even if those mechanisms are perfect.

\paragraph{Stop point.}
This note stops once a deployment has a declared committee-contact witness and the corresponding MUCC contact-floor inequalities. At that point, the committee-contact claim is referee-complete. Any richer endpoint semantics, evaluator notarization, runtime monitoring, or release packaging belong to the companion anonymity/state/evaluation/release notes already cited above.\cite{series:oddsinflation,series:profiling,series:state3,series:evalnotary,series:release}


\section{Model and main results}
\label{sec:model}

\paragraph{Committees and destinations.}

\paragraph{Witness projection (CCT $\rightarrow$ committee contact).}
We treat a lookup as producing a committee-contact trace (CCT) as in Synthesis~3.
MUCC concerns a specific projection of that trace: the indicator (or set) of committees contacted during the lookup.
Formally, fix an exposure normal form \(\mathsf{ENF}\) whose \texttt{exposure\_nf\_id} names the projection \(g(\mathrm{CCT})\); all probabilities below are with respect to that declared \(g\).

A lookup targets a key whose responsible set is a destination committee-set $D\subseteq[M]$ of size $|D|=d$ (replication factor).
A lookup succeeds if it obtains at least $t$ \emph{live} responses from committees in $D$.
Let $\alpha\in(0,1]$ denote the probability that a contacted destination committee yields a live response (capturing churn/drops and adversarial suppression).
Section~\ref{sec:churn_caching} gives a simple mapping from churn/caching models into a conservative lower bound on $\alpha$.

\paragraph{MUCC marginals.}
Let $C$ be the (random) set of committees contacted over the \emph{entire} lookup (possibly across multiple rounds/subqueries).
For multi-round protocols we write $C^{(\ell)}$ for the round-$\ell$ contact set, so that $C=\bigcup_{\ell=1}^L C^{(\ell)}$.

\begin{definition}[MUCC (marginally-uniform committee contact) marginals]
\label{def:mucc}
A lookup mechanism has \emph{MUCC marginals} if for every key (hence every destination set $D$) and every committee $j\in[M]$,
\[\Pr[j\in C\mid D] = p\quad\text{(a constant independent of $D$ and $j$)}.\]
Write $k_0 := \E[|C|] = Mp$.
\end{definition}

\begin{remark}[MUCC is marginal-only]
MUCC constrains only the \emph{single-committee} marginals $\Pr[j\in C\mid D]$.
Equivalently, for each committee $j$, the one-bit projection $X_j:=\one\{j\in C\}$ is marginally equalized across destinations.
The \emph{joint} distribution of the contacted set $C$ may still depend on $D$ and leak the key.

Accordingly, MUCC is best viewed as a \emph{first-moment footprint contract}:
it is weaker than the full equalization contracts used for profiling budgets, which bound the \emph{distribution} of a declared transcript projection $X=f(T)$ (Anonymity~B).\cite{series:profiling}
Our ceilings/floors below therefore use MUCC only as a \emph{minimal} footprint-hiding assumption:
they apply \emph{a fortiori} to any stronger contact-pattern privacy notion (e.g., full key-independence of $C$, or DP-style stability over contact sets).
\end{remark}

MUCC captures the ``key-independent footprint'' objective: the observer cannot learn the destination key from single-committee marginals.

\begin{lemma}[Destination contact expectation under MUCC]
\label{lem:mucc-dest-exp}
Let $Y:=|C\cap D|$ be the number of destination committees contacted over the whole lookup.
Under MUCC, $\E[Y]=dp$.
If each contacted destination committee yields a live response with probability $\alpha$ (independent thinning), then the number of live destination responses $Z$ satisfies $\E[Z]=\alpha dp$.
\end{lemma}

\begin{proof}
By linearity of expectation, $\E[Y]=\sum_{i\in D}\Pr[i\in C\mid D]=dp$.
Under independent thinning with rate $\alpha$, $\E[Z]=\alpha\E[Y]=\alpha dp$.
\end{proof}

\subsection{Availability ceilings under MUCC}
By Lemma~\ref{lem:mucc-dest-exp}, MUCC fixes the expected number of live destination responses: $\E[Z]=\alpha dp$.
This is the only property we need for the universal ceiling below.

\begin{theorem}[MUCC availability ceiling (universal, expectation form)]
\label{thm:mucc-ceiling}
Consider any lookup mechanism with MUCC marginals.
If the lookup succeeds with probability at least $1-\beta$ (i.e., $\Pr[Z\ge t]\ge 1-\beta$), then necessarily
\[
 d\,p\,\alpha \;\ge\; (1-\beta)\,t.
\]
Equivalently, the expected total contact satisfies
\[
 k_0 \;\ge\; \frac{M\,(1-\beta)\,t}{\alpha\,d}.
\]
\end{theorem}

\begin{proof}[Proof sketch]
For nonnegative $Z$, Markov gives $\Pr[Z\ge t]\le \E[Z]/t$.
Rearrange.
\end{proof}

\begin{remark}[Union-only observers]
If the observer sees only the union footprint $C$ (not per-round contacts), then Theorem~\ref{thm:mucc-ceiling} already lower bounds $\E[|C|]=Mp$.
Since total contact events satisfy $K\ge |C|$ always, the same floor applies to $\E[K]$.
Proposition~\ref{prop:mucc-stop-time-floor} isolates the common (and stricter) regime where per-round footprints are observable.
\end{remark}

\begin{remark}[Tightness of the expectation ceiling]
Theorem~\ref{thm:mucc-ceiling} is proved via Markov's inequality and is essentially tight without additional assumptions.
Indeed, for any $\mu\in[0,t]$ there exists a nonnegative random variable $Z$ with $\E[Z]=\mu$ and $\Pr[Z\ge t]=\mu/t$ (take $Z=t$ with probability $\mu/t$ and $Z=0$ otherwise).
Thus, improving Theorem~\ref{thm:mucc-ceiling} requires adding structure (e.g., independence, sampling without replacement, or moment bounds).
\end{remark}

\begin{corollary}[Approximate MUCC ceilings (additive slack)]
\label{cor:approx-mucc}
Fix $\eta\in[0,1]$.
Suppose a lookup mechanism satisfies an \emph{$\eta$-approximate MUCC upper envelope}:
there exists a constant $p$ (independent of $D$ and $j$) such that for every destination set $D$ and every committee $j\in[M]$,
\[
\Pr[j\in C\mid D] \le p+\eta.
\]
If the lookup succeeds with probability at least $1-\beta$, then necessarily
\[
p \;\ge\; \frac{(1-\beta)\,t}{\alpha\,d}-\eta,
\qquad\text{hence}\qquad
Mp \;\ge\; \frac{M\,(1-\beta)\,t}{\alpha\,d}-M\eta.
\]
In other words, relaxing MUCC by $\eta$ can reduce the MUCC ceiling by at most an additive $M\eta$ term.
\end{corollary}

\begin{proof}[Proof sketch]
Under the envelope, $\E[Z]\le d\,(p+\eta)\alpha$.
Combine with Markov as in Theorem~\ref{thm:mucc-ceiling}.
\end{proof}


The bound is intentionally assumption-light (no independence required): it uses only MUCC marginals and a live-response factor $\alpha$.
Sharper ceilings are possible under additional structure (e.g., uniform sampling without replacement) using hypergeometric concentration inequalities~\cite{serfling1974sampling,bardenet2015withoutreplacement}, but Theorem~\ref{thm:mucc-ceiling} is already strong in large-$M$ regimes.


\subsection{A best-case tail ceiling under i.i.d.\ destination success}
The expectation ceiling above is universal and assumption-light, but it does not capture the additional slack needed when $d$ is moderate and one wants small failure probability.
To state a clean strengthening, assume an \emph{optimistic} model in which for each destination committee $i\in D$,
the event ``$i$ yields a live response'' is independent with probability $p\alpha$ (e.g., contact decisions are approximately independent across $i$ and liveness is independent thinning).

\begin{proposition}[Binomial-quantile ceiling (optimistic i.i.d.\ model)]
\label{prop:binom-quantile-ceiling}
Under the i.i.d.\ destination-success model, $Z\sim \mathrm{Binomial}(d,p\alpha)$.
Define
\[
p^\star(d,t,\beta) := \inf\{q\in[0,1] : \Pr[\mathrm{Binomial}(d,q)\ge t]\ge 1-\beta\}.
\]
If the lookup succeeds with probability at least $1-\beta$, then necessarily
\[
p\alpha \;\ge\; p^\star(d,t,\beta),
\qquad\text{hence}\qquad
k_0 = Mp \;\ge\; \frac{M}{\alpha}\,p^\star(d,t,\beta).
\]
\end{proposition}

\begin{proof}
Under the i.i.d. model, $Z\sim\mathrm{Binomial}(d,p\alpha)$.
The success requirement $\Pr[Z\ge t]\ge 1-\beta$ implies that the Binomial success probability parameter $p\alpha$ must be at least the minimal value $p^\star(d,t,\beta)$ achieving that tail probability.
Rearrange to obtain the stated lower bound on $k_0=Mp$.
\end{proof}

\begin{lemma}[A closed-form conservative ceiling (Hoeffding)]
\label{lem:hoeffding-ceiling}
Under the i.i.d.\ destination-success model, if
\[
p\alpha \;\ge\; \frac{t}{d} + \sqrt{\frac{\log_2(1/\beta)}{2d}},
\]
then $\Pr[Z\ge t]\ge 1-\beta$.
Equivalently, a sufficient MUCC contact fraction is
\[
\frac{k_0}{M}=p \;\ge\; \frac{1}{\alpha}\left(\frac{t}{d} + \sqrt{\frac{\log_2(1/\beta)}{2d}}\right).
\]
\end{lemma}

\begin{proof}[Proof sketch]
Let $Z\sim\mathrm{Binomial}(d,p\alpha)$ and write $q=p\alpha$.
Hoeffding's inequality gives $\Pr[Z/d \le q-\epsilon]\le \exp(-2d\epsilon^2)$.
Setting $\epsilon = q - t/d$ and requiring $\exp(-2d\epsilon^2)\le\beta$ yields the condition.
\end{proof}

Table~\ref{tab:binom-quantile} gives representative values; correlations and active adversaries can only \emph{increase} the contact required to reach a given success probability.

\subsection{A best-case ceiling under uniform fixed-size contact (hypergeometric)}
MUCC alone does not specify how the contact set $C$ is sampled.
To sharpen ceilings under a concrete (optimistic) sampling model, assume $C$ is sampled as a \emph{uniform random subset of fixed size} $|C|=k_0$, independent of $D$.
This is a best-case MUCC instantiation: it equalizes marginals and minimizes variance among exchangeable schemes.

Under this model, $Y:=|C\cap D|$ follows a hypergeometric law $Y\sim \mathrm{Hypergeom}(M,d,k_0)$ with mean $\E[Y]=dk_0/M$.
If each contacted destination committee yields at most one response, and yields a live response with probability at least $\alpha$, then a conservative necessary condition for obtaining $t$ live responses is that $Y\ge t'$, where $t':=\lceil t/\alpha\rceil$.
Therefore a necessary ceiling is:
\[
\Pr\!\bigl[\mathrm{Hypergeom}(M,d,k_0)\ge t'\bigr]\ \ge\ 1-\beta.
\]
Equivalently, define the hypergeometric quantile
\[
k_0^{\star}(M,d,t',\beta)\ :=\ \inf\{k\in\{0,\dots,M\}:\Pr[\mathrm{Hypergeom}(M,d,k)\ge t']\ge 1-\beta\},
\]
and note that any uniform fixed-size MUCC scheme meeting the success target must satisfy $k_0\ge k_0^\star$.
Closed-form conservative approximations are available via Serfling/without-replacement inequalities.\cite{serfling1974sampling,bardenet2015withoutreplacement}

\subsection{Abort-safe equalization implies a total cover-contact floor}
To hide retries/early success, anonymous lookups often fix a stop time (number of rounds/subqueries) and run dummy work when success would otherwise happen early.
Let $L$ be the fixed number of rounds, and let $C^{(\ell)}$ be the contact set in round $\ell$.

\begin{proposition}[Abort-safe cover-contact floor (splitting into rounds does not help)]
\label{prop:mucc-stop-time-floor}
Assume the observer sees per-round contacts and that each round has MUCC marginals with a common per-round contact probability $p_{\mathrm{rnd}}$ (so $\E[|C^{(\ell)}|]=Mp_{\mathrm{rnd}}$).
Let $K:=\sum_{\ell=1}^L |C^{(\ell)}|$ be the total number of committee-contact events executed under abort-safe equalization (including dummies).
If the lookup succeeds with probability at least $1-\beta$, then necessarily
\[
\E[K] \;\ge\; \frac{M\,(1-\beta)\,t}{\alpha\,d}.
\]
Equivalently, more rounds can reduce per-round contact ($p_{\mathrm{rnd}}\propto 1/L$ in the best case), but cannot reduce the required \emph{total} cover-contact.
\end{proposition}

\begin{proof}[Proof sketch]
Under per-round MUCC, the expected number of live destination responses across all rounds is
$\E[Z]=\alpha\,d\,\sum_{\ell=1}^L \Pr[i\in C^{(\ell)}\mid D] = \alpha\,d\,L p_{\mathrm{rnd}}$.
Apply Markov as in Theorem~\ref{thm:mucc-ceiling}: $\Pr[Z\ge t]\le \E[Z]/t$.
Finally, $\E[K]=\sum_{\ell=1}^L \E[|C^{(\ell)}|]=L M p_{\mathrm{rnd}}$.
\end{proof}

\begin{figure}[t]
  \centering
  \fbox{\parbox{0.92\linewidth}{\centering\textbf{Figure omitted in source-only archive.}\\(Plot/diagram and scripts available on request.)}}
  \caption{Illustrative abort-safe equalization tradeoff: hiding retries/early success imposes an unavoidable ``run the full schedule'' cost.}
  \label{fig:abort-eq}
\end{figure}


\section{Implications and design notes}
\label{sec:implications}

\paragraph{What the ceiling means.}
Theorem~\ref{thm:mucc-ceiling} says that in large deployments, hiding destination from contact-set marginals can force a lookup to ``touch'' a substantial fraction of the network when availability targets are strict.
In particular, the ceiling worsens as $\alpha$ falls (churn/drops) or as $t$ rises (stronger reconstruction/availability).

\paragraph{Worked regime table.}
To keep this addendum referee-tight but brief, we include two TeX-only evidence artifacts:
(i) a small calibration table translating $(M,d,t,\alpha,\beta)$ into the implied $k_0$ floor,
and (ii) an MLDHT scale note motivating the large-$M$ regime.

\begin{table}[h]
\centering
\resizebox{\linewidth}{!}{%
\begin{tabular}{@{}rrrr@{}}
\toprule
Live factor $\alpha$ & Threshold $t$ (of $d=8$) & Required contact fraction $k_0/M$ & Feasibility \\
\midrule
1.0 & 3 & 0.356 & feasible \\
1.0 & 4 & 0.475 & feasible \\
1.0 & 5 & 0.594 & feasible \\
0.8 & 3 & 0.445 & feasible \\
0.8 & 4 & 0.594 & feasible \\
0.8 & 5 & 0.742 & feasible \\
0.6 & 3 & 0.594 & feasible \\
0.6 & 4 & 0.792 & feasible \\
0.6 & 5 & 0.990 & \emph{near-total contact} \\
0.4 & 3 & 0.891 & feasible (high overhead) \\
0.4 & 4 & 1.187 & \emph{impossible under MUCC} \\
0.4 & 5 & 1.484 & \emph{impossible under MUCC} \\
\bottomrule
\end{tabular}%
}
\caption{Illustrative lower bounds from Theorem~\ref{thm:mucc-ceiling} with $\beta=0.05$.
The required contact fraction depends mainly on $t/(\alpha d)$.
Values above $1$ indicate that no MUCC-marginal mechanism can meet the success target under the stated live factor.
}
\label{tab:mucc-numbers}
\end{table}


\begin{table}[h]
\centering
\begin{tabular}{@{}rrr@{}}
\toprule
Live membership $p_{\mathrm{live}}$ & Drop rate $p_{\mathrm{drop}}$ & Effective $\alpha=p_{\mathrm{live}}(1-p_{\mathrm{drop}})$ \\
\midrule
0.9 & 0.05 & 0.855 \\
0.8 & 0.10 & 0.720 \\
0.7 & 0.15 & 0.595 \\
0.6 & 0.20 & 0.480 \\
\bottomrule
\end{tabular}
\caption{A simple factorization of the live-response probability $\alpha$ into membership and per-contact drop/suppression.
This table is illustrative and intended only for interpreting the ceiling in Theorem~\ref{thm:mucc-ceiling}.}
\label{tab:alpha}
\end{table}


\begin{table}[h]
\centering
\begin{tabular}{@{}rrrrr@{}}
\toprule
$d$ & $t$ & $p^\star$ for $1-\beta=0.95$ & $p^\star$ for $1-\beta=0.99$ & $p^\star$ for $1-\beta=0.999$ \\
\midrule
20 & 14 & 0.823 & 0.871 & 0.914 \\
32 & 23 & 0.820 & 0.860 & 0.897 \\
64 & 45 & 0.782 & 0.815 & 0.847 \\
\bottomrule
\end{tabular}
\caption{Exact binomial-quantile ceilings under the optimistic i.i.d.\ destination-success model $Z\sim\mathrm{Binomial}(d,p\alpha)$: 
$p^\star(d,t,\beta)$ is the smallest success rate such that $\Pr[Z\ge t]\ge 1-\beta$. 
To translate to MUCC, use $p = k_0/M$ and thus $k_0 \ge (M/\alpha)\,p^\star$. 
Rows use thresholds $t\approx 0.7d$ to illustrate that \emph{high} success probability forces $p\alpha$ well above $t/d$ even before considering correlations/adversaries.}
\label{tab:binom-quantile}
\end{table}


\paragraph{MLDHT scale note (evidence).}
Mainline BitTorrent's DHT is specified by BEP-5~\cite{bep5}.
Measurement work reports deployments at the scale of tens of millions of nodes with substantial churn~\cite{wang2013mldht,mldhtProjectPage}.
As a complementary ``load/usage'' indicator, MagnetDB reports 28.6M torrents discovered via the DHT over 2018--2024~\cite{magnetdb2025}.
The ceiling in Theorem~\ref{thm:mucc-ceiling} is most constraining in exactly these large-$M$ regimes.



\subsection{Churn and caching as \texorpdfstring{$\alpha$}{alpha}}
\label{sec:churn_caching}
Theorems~\ref{thm:mucc-ceiling} and~\ref{prop:binom-quantile-ceiling} treat $\alpha$ as an external ``live response'' factor.
To make the ceiling plug-and-play, it is useful to map protocol assumptions (uptime, stabilization delay, cache staleness) into a conservative lower bound on $\alpha$.
We state one clean mapping under standard independence approximations.

\begin{lemma}[A simple $\alpha$ bound from churn and quorum]
\label{lem:alpha-from-churn}
Suppose a destination committee consists of $d$ nodes and a response is ``live'' if at least $t$ of the $d$ nodes are reachable.
Assume (i) each node is reachable at the moment it is queried with probability at least $q$ (independent across nodes), and (ii) reachability already includes any fixed stabilization delay $\Delta$ and any cache-staleness penalty.
Then the committee live-response probability satisfies
\[
\alpha\;\ge\;\Pr[\mathrm{Binomial}(d,q)\ge t].
\]
If node sessions have survival function $S(\Delta)=\Pr[\text{session length}\ge \Delta]$ and per-node query success probability at zero delay is at least $q_0$, then one conservative choice is $q=q_0\,S(\Delta)$.
\end{lemma}

\begin{proof}[Proof sketch]
Under the assumptions, the number of reachable nodes stochastically dominates a $\mathrm{Binomial}(d,q)$ random variable.
\end{proof}

\begin{remark}[A one-line correlated-outage refinement]
If, with probability $r$, the destination committee experiences a common-mode outage (e.g., shared ASN/rack failure) that makes all $d$ nodes unreachable, and otherwise nodes are independently reachable with probability at least $q$, then
\[\alpha\;\ge\;(1-r)\,\Pr[\mathrm{Binomial}(d,q)\ge t].\]
This simple mixture model is often enough to see how quickly the MUCC ceiling worsens under correlated failures.
\end{remark}

Lemma~\ref{lem:alpha-from-churn} is intentionally coarse: correlations, adversarial suppression, and load-dependent failure can only \emph{decrease} $\alpha$.
The evidence table \texttt{churn\_alpha\_table.tex} provides representative values to help calibrate how quickly the MUCC ceiling worsens as $\alpha$ drops.

\paragraph{Realizing MUCC marginals.}
A canonical MUCC scheduler chooses $k_0$ committees uniformly without replacement each round.
To maintain state across rounds while preserving uniform marginals, one can evolve a fixed-size contact set via a random walk on the Johnson graph (Bernoulli--Laplace diffusion), which mixes rapidly in $k_0$ and provides a simple implementation path~\cite{diaconis1987bernoulliLaplace}.

\paragraph{Relationship to other series results.}
Schedule equalization and stop-time padding (from the series) can remove additional witness channels (timing, retries), but they cannot evade the MUCC ceiling once MUCC is required.
Paper~1 provides the broader state-dependent lens (amplification and accounting) for combining these constraints.

\begin{figure}[t]
  \centering
  \fbox{\parbox{0.92\linewidth}{\centering\textbf{Figure omitted in source-only archive.}\\(Plot/diagram and scripts available on request.)}}
  \caption{Illustrative boundary between cover-contact schedules.}
\end{figure}


\section{Related work}
Structured overlays and DHTs include Chord, Pastry, and Kademlia~\cite{stoica2001chord,rowstron2001pastry,maymounkov2002kademlia}.
Anonymous or privacy-enhanced overlays built on structured topologies include Salsa, ShadowWalker, and AP3~\cite{nambiar2006salsa,mittal2009shadowwalker,mislove2006ap3}.
Attacks on anonymous DHT-style lookups (and the motivation for secure+anonymous lookup) were developed by Wang et al.~\cite{wang2010anonymousLookupAttacks,mittal2012informationleaks}.
Routing-table compromise as an anonymity/metadata leak channel (for storage/recipient anonymity) was quantified in an information-theoretic framework by Ray and Zhang~\cite{ray2007itframework}.
Octopus aims to provide both security and anonymity for DHT lookup via splitting and attacker identification~\cite{wang2012octopus}.

On the ``query privacy'' axis, Backes et al.~\cite{backes2012queryprivacy} study adding query privacy to robust DHTs.
More recently, Daniel et al.~\cite{daniel2025kademliaObfuscation} propose query-obfuscation methods for Kademlia-style lookups.
Aktypi and Rasmussen's Iris system targets query privacy in authenticated Chord networks while remaining compatible with non-upgraded peers~\cite{aktypi2025iris}.
These works improve \emph{content} privacy of the queried key, but (unless combined with additional mechanisms) still leave a contact-footprint witness channel: which peers/committees were contacted and in what pattern.
For deployed-scale context, BEP-5 specifies the Mainline BitTorrent DHT~\cite{bep5}.
Large-scale measurement studies highlight that Mainline DHT deployments can involve tens of millions of nodes with substantial churn~\cite{wang2013mldht,mldhtProjectPage}.
BitTorrent DHT crawling and measurement attacks further motivate treating contact footprints as a realistic observation channel~\cite{wolchok2010crawling,falkner2007profiling}.
Real-world Sybil attacks have also been observed at scale in Mainline DHT deployments, reinforcing that adversaries can both observe and manipulate contact footprints~\cite{wang2012sybil}.

Cryptographic \emph{query-privacy} mechanisms (PIR/PSI-style) can hide the lookup item from contacted peers, but they do not by themselves hide \emph{which} peers/committees are contacted; the contact-set witness channel remains.

Beyond IPFS, I2P's netDB uses a DHT-like floodfill subset for storing and locating RouterInfos and LeaseSets; entries are replicated to the three floodfills closest to a routing key and lookups are often sent in parallel to two closest ``good'' floodfills~\cite{i2p_netdb_overview}.
Notably, the operational docs also distinguish confidentiality across object types: LeaseSet lookups are end-to-end garlic encrypted (hiding content from intermediate routers), while RouterInfo lookups are not encrypted and can be snooped by the outbound endpoint of the client tunnel~\cite{i2p_netdb_overview}.

Finally, IPFS uses a Kademlia-derived DHT for peer/content discovery~\cite{ipfs_kad_dht_spec}.
This surface has been shown vulnerable to Sybil/eclipse and censorship behaviors, including provider-record censorship in NDSS'24~\cite{sridhar2024ipfs_censorship}, single-machine Sybil denial of content access~\cite{cholez2024ipfs_sybil}, follow-on active Sybil refinements~\cite{netto2025ipfs_active_sybil}, and network-level/BGP censorship models~\cite{mattertran2025ipfs_censorship}.
On the privacy side, Peer2PIR integrates PIR into IPFS routing to reduce query-content leakage~\cite{mazmudar2024peer2pir}, while Karapapas et al. document anonymity abuse patterns and gateway/pinning behavior in the wild~\cite{karapapas2025ipfs_anonymity_abuse}.
Our MUCC analysis is complementary: it isolates the availability/contact-pattern floor that persists even when query contents are obfuscated.




\section{Protocol mapping: plugging a design into \texorpdfstring{$(M,d,t,\alpha,k_0)$}{(M,d,t,alpha,k0)}}\label{sec:protocol-mapping}

This addendum is stated at the level of a destination committee-set $D\subseteq[M]$ (size $d$) and a required live-response threshold $t$.
To apply it to a concrete anonymous DHT or directory design, extract the following from the protocol specification:
(i) what constitutes a ``destination committee'' for a key (often a replica group or a set of responsible nodes),
(ii) how many such committees exist system-wide ($M$),
(iii) how many are responsible for a key ($d$) and what success condition is required ($t$),
(iv) the committee live-response probability $\alpha$ (Section~\ref{sec:churn_caching}), and
(v) the expected total distinct committees contacted $k_0=\E[|C|]$ (including dummies and retries under abort-safe equalization).

\begin{table}[t]
\centering
\small
\begin{tabularx}{\linewidth}{@{}>{\raggedright\arraybackslash}p{0.27\linewidth} Y Y@{}}
\toprule
\textbf{Design pattern} & \textbf{How to read $(M,d,t)$} & \textbf{Notes on $k_0$ and $\alpha$} \\
\midrule
\textbf{Octopus-style anonymous lookup}~\cite{wang2012octopus}
& Treat each DHT node (or bucket) as a ``committee'': $M\approx N$, $d=1$, $t=1$.
& $k_0$ counts all queried nodes across split paths + dummy queries; $\alpha$ is per-target response probability (path reliability + churn). \\

\textbf{Salsa / ShadowWalker / NISAN-style directories}~\cite{nambiar2006salsa,mittal2009shadowwalker,nisan2009}
& ``Destination'' is a directory entry / successor set; typically $d>1$ (redundant lookups) with $t\ge 1$.
& Redundant lookups and bounds checking increase $k_0$; churn and cache staleness reduce $\alpha$ and therefore tighten the MUCC ceiling. \\

\textbf{Replicated DHT with $d$ replicas} (generic)
& $M$ is the number of replica groups/buckets; $d$ is replicas per key; $t=1$ for ``any replica answers''.
& Even if queries are PIR-obfuscated (e.g., Peer2PIR~\cite{mazmudar2024peer2pir}), MUCC concerns contact-set marginals: $k_0$ must still cover enough groups to make footprints key-independent. \\

\textbf{IPFS {K}ademlia DHT provider records}~\cite{ipfs_kad_dht_spec}
& A content key is stored/served by the $k$ closest peers ($k=20$ in the IPFS spec), so a natural mapping is $d=k$ and $t=1$ for ``any answer'' (or $t>1$ for multi-source confirmation).
& Provider records expire after $48$h and are recommended to be republished every $22$h in the IPFS spec~\cite{ipfs_kad_dht_spec}; expiry/churn reduces $\alpha$ and tightens MUCC ceilings. \\

\textbf{I2P netDB (floodfill DHT subset)}~\cite{i2p_netdb_overview}
& Treat floodfill routers as the committee population: $M\approx$ \#floodfills. NetDB entries are ``flooded'' to the \emph{three} floodfills closest (by XOR distance) to a routing key, so a conservative mapping is $d=3$ and $t=1$ for ``any answer'' (or $t>1$ if clients require multi-floodfill confirmation).
& The docs note lookups are issued in parallel to the \emph{two} closest ``good'' floodfills; LeaseSet lookups are end-to-end garlic encrypted while RouterInfo lookups are not encrypted (and are snoopable by the outbound endpoint), so contact footprints remain a natural witness stream. Redundancy and verification attempts increase $k_0$ and raise the MUCC floor.\\


\textbf{Threshold / secret-shared destinations} (generic)
& $t>1$ (quorum/threshold reconstruction) makes the ceiling stronger: $k_0\ge M(1-\beta)t/(\alpha d)$.
& Abort-safe equalization forces the full schedule: splitting into $L$ rounds does \emph{not} reduce the required total cover-contact (Proposition~\ref{prop:mucc-stop-time-floor}). \\
\bottomrule
\end{tabularx}
\caption{One-page mapping from common anonymous DHT/directory design patterns into the parameters of Theorem~\ref{thm:mucc-ceiling}.
The mapping is intentionally conservative and ignores correlations that would only worsen the ceiling.
}
\label{tab:protocol-mapping}
\end{table}





\section{Takeaways}
\begin{itemize}[leftmargin=*]
\item \textbf{Footprints have floors:} MUCC marginals force nontrivial cover-contact; success and availability targets translate into a required \emph{fraction of all committees} contacted.
\item \textbf{Abort-safe equalization does not help:} if stop-time is fixed to hide retries/aborts, splitting into more rounds cannot reduce the total cover-contact implied by MUCC.
\item \textbf{System parameters matter:} higher per-committee liveness $\alpha$ and larger replica sets $d$ loosen MUCC ceilings; churn/caching correlations tighten them.
\item \textbf{This is the committee-contact root:} once the MUCC packet is fixed, profiling, monitoring, notarization, and release papers should cite it rather than restating the ceiling argument.
\item \textbf{Use MUCC deliberately:} MUCC is a strong contract surface; if you only need to hide timing or hop patterns, equalize those projections instead (Anonymity~B) rather than forcing MUCC everywhere.
\end{itemize}


\begin{thebibliography}{99}\small

\bibitem{series:schedcompiler}
Anonymous.
\newblock \emph{Scheduling and Compiler Techniques for Metadata-Hiding Overlay Lookups}.
\newblock Series paper (published), 2026.
\newblock Wiki: \texttt{[[2026.01.22 - Mathematics: Scheduling and Compiler Techniques for Metadata-Hiding Overlay Lookups]]}.

\bibitem{series:stoptimeaddendum}
Anonymous.
\newblock \emph{Stop-Time Padding Addendum (unique-only): Max-Mixture Separations and Tail-Sign Deadline Normal Forms}.
\newblock Series note (published), 2026.
\newblock Wiki: \texttt{[[2026.01.26 - Mathematics: Stop-Time Padding Addendum, Max-Mixture Separations and Tail-Sign Deadline Normal Forms]]}.

\bibitem{series:operationalA}
Anonymous.
\newblock Schedule-Only Retries in Anonymous DHTs: SRP and Retry-Trace Equalization.
\newblock Operational series Paper~A, 2026. (This archive.)

\bibitem{series:operationalB}
Anonymous.
\newblock Auditable Trace Privacy with Abort: Plan-Pinning and Drift Receipts for Side-Channel Resistant Transport.
\newblock Operational series Paper~B, 2026. (This archive.)

\bibitem{series:state1}
Anonymous.
\newblock Stateful Anonymity and Control-Plane Leakage: Selection Privacy, Amplification, and Accounting.
\newblock Companion preprint in this archive (AnonDHT~State~1).

\bibitem{series:profiling}
Anonymous.
\newblock Anonymous DHT Profiling: Trace Projections, Equalization Contracts, and Identification Caps.
\newblock Companion preprint in this archive (Anonymity~B).

\bibitem{series:state3}
Anonymous.
\newblock Closed-View Auditing for Stateful Anonymity: Odometers, Alarms, and Machine-Checkable CPPCs.
\newblock Companion preprint in this archive (AnonDHT~State~3).

\bibitem{nambiar2006salsa}
Arjun Nambiar and Matthew Wright.
\newblock Salsa: A Structured Approach to Large-Scale Anonymity.
\newblock Proceedings of the 13th ACM Conference on Computer and Communications Security (CCS), 17--26, 2006.
\newblock DOI: 10.1145/1180405.1180409.

\bibitem{mittal2009shadowwalker}
Prateek Mittal and Nikita Borisov.
\newblock ShadowWalker: Peer-to-peer Anonymous Communication Using Redundant Structured Topologies.
\newblock ACM Conference on Computer and Communications Security (CCS), 161--172, 2009.
\newblock DOI: 10.1145/1653662.1653683.
\newblock CCS 2009; open-access mirror via anonbib/freehaven.
\newblock \url{https://www.freehaven.net/anonbib/cache/ccs09-shadowwalker.pdf}

\bibitem{wang2012octopus}
Qiyan Wang and Nikita Borisov.
\newblock Octopus: A Secure and Anonymous {DHT} Lookup.
\newblock IEEE International Conference on Distributed Computing Systems (ICDCS), 325--334, 2012.
\newblock DOI: 10.1109/ICDCS.2012.78.
\newblock Also available as arXiv:1203.2668; open author PDF.
\newblock \url{https://nikita.ca/papers/octopus-icdcs12.pdf}

\bibitem{wang2010anonymousLookupAttacks}
Qiyan Wang and Prateek Mittal and Nikita Borisov.
\newblock In Search of an Anonymous and Secure Lookup: Attacks on Structured Peer-to-Peer Anonymous Communication Systems.
\newblock ACM Conference on Computer and Communications Security (CCS), 308--318, 2010.
\newblock DOI: 10.1145/1866307.1866343.
\newblock CCS 2010; open-access mirror via anonbib/freehaven.
\newblock \url{https://www.freehaven.net/anonbib/cache/ccs10-lookup.pdf}

\bibitem{nisan2009}
Andriy Panchenko and Stefan Richter and Arne Rache.
\newblock {NISAN}: Network Information Service for Anonymization Networks.
\newblock Proceedings of the 16th {ACM} Conference on Computer and Communications Security (CCS '09), 141--150, 2009.
\newblock DOI: 10.1145/1653662.1653681.
\newblock DHT-based network-information distribution designed to mitigate lookup bias/leaks in anonymization networks.

\bibitem{mittal2012informationleaks}
Prateek Mittal and Nikita Borisov.
\newblock Information Leaks in Structured Peer-to-Peer Anonymous Communication Systems.
\newblock ACM Transactions on Information and System Security, 15(1), 1--28, 2012.
\newblock DOI: 10.1145/2133375.2133380.
\newblock Article 5 (March 2012), 28 pages; open author PDF.
\newblock \url{https://www.princeton.edu/~pmittal/publications/information-leaks-tissec.pdf}

\bibitem{ray2007itframework}
Souvik Ray and Zhao Zhang.
\newblock An Information-Theoretic Framework for Analyzing Leak of Privacy in Distributed Hash Tables.
\newblock Proceedings of the 7th {IEEE} International Conference on Peer-to-Peer Computing (P2P), 27--36, 2007.
\newblock Quantifies privacy loss from compromised routing tables; open author PDF.
\newblock \url{https://home.engineering.iastate.edu/zzhang/publications/tr07-p2p-privacy.pdf}

\bibitem{stoica2001chord}
Ion Stoica and Robert Morris and David Karger and M. Frans Kaashoek and Hari Balakrishnan.
\newblock Chord: A Scalable Peer-to-peer Lookup Service for Internet Applications.
\newblock ACM SIGCOMM, 2001.

\bibitem{rowstron2001pastry}
Antony Rowstron and Peter Druschel.
\newblock Pastry: Scalable, Decentralized Object Location, and Routing for Large-Scale Peer-to-Peer Systems.
\newblock IFIP/ACM International Conference on Distributed Systems Platforms (Middleware), 2001.

\bibitem{maymounkov2002kademlia}
Petar Maymounkov and David Mazi{\`e}res.
\newblock Kademlia: A Peer-to-peer Information System Based on the XOR Metric.
\newblock International Workshop on Peer-to-Peer Systems (IPTPS), 2002.

\bibitem{mislove2006ap3}
Alan Mislove and Atul Oberoi and Ansley Post and Peter Druschel and Dan S. Wallach.
\newblock AP3: Cooperative, Decentralized Anonymous Communication.
\newblock ACM SIGOPS/EuroSys European Conference on Computer Systems, 2006.

\bibitem{backes2012queryprivacy}
Michael Backes and Ian Goldberg and Aniket Kate and Tomas Toft.
\newblock Adding Query Privacy to Robust DHTs.
\newblock Proceedings of the 7th ACM Symposium on Information, Computer and Communications Security (ASIACCS), 30--31, 2012.
\newblock DOI: 10.1145/2414456.2414473.
\newblock Extended version: arXiv:1107.1072.

\bibitem{daniel2025kademliaObfuscation}
Erik Daniel and Guillaume Michel and Florian Tschorsch.
\newblock Improving {Kademlia} Lookup Privacy through Query Obfuscation.
\newblock Proceedings of the 40th {ACM}/{SIGAPP} Symposium on Applied Computing (SAC), 1345--1352, 2025.
\newblock DOI: 10.1145/3672608.3707703.
\newblock \url{https://dl.acm.org/doi/10.1145/3672608.3707703}

\bibitem{aktypi2025iris}
Angeliki Aktypi and Kasper Rasmussen.
\newblock Iris: Dynamic Privacy Preserving Search in Authenticated Chord Peer-to-Peer Networks.
\newblock Proceedings of the 32nd Annual Network and Distributed System Security Symposium (NDSS), 2025.
\newblock DOI: 10.14722/ndss.2025.240392.
\newblock \url{https://www.ndss-symposium.org/ndss-paper/iris-dynamic-privacy-preserving-search-in-authenticated-chord-peer-to-peer-networks/}

\bibitem{diaconis1987bernoulliLaplace}
Diaconis, Persi and Shahshahani, Mehrdad.
\newblock Time to reach stationarity in the {B}ernoulli--{L}aplace diffusion model.
\newblock {SIAM} Journal on Mathematical Analysis, 18(1), 208--218, 1987.
\newblock DOI: 10.1137/0518016.

\bibitem{falkner2007profiling}
Jarret Falkner and Michael Piatek and John P. John and Arvind Krishnamurthy and Thomas E. Anderson.
\newblock Profiling a million user DHT.
\newblock Proceedings of the 7th ACM SIGCOMM Conference on Internet Measurement (IMC), 129--134, 2007.
\newblock DOI: 10.1145/1298306.1298325.

\bibitem{cortes2024kademlia}
Mikel Cortes-Goicoechea and Csaba Kiraly and Dmitriy Ryajov and Jose Luis Mu{\~n}oz-Tapia and Leonardo Bautista-Gomez.
\newblock Scalability limitations of Kademlia DHTs when enabling Data Availability Sampling in Ethereum.
\newblock 2024.
\newblock arXiv:2402.09993

\bibitem{wolchok2010crawling}
Scott Wolchok and J. Alex Halderman.
\newblock Crawling {BitTorrent} {DHTs} for Fun and Profit.
\newblock 4th USENIX Workshop on Offensive Technologies (WOOT 10), 2010.
\newblock Open-access PDF available via USENIX (Wolchok.pdf).
\newblock \url{https://www.usenix.org/conference/woot10/crawling-bittorrent-dhts-fun-and-profit}

\bibitem{ourManyTesting}
Anonymous Series Authors.
\newblock Many-Testing Lower Bounds for Hierarchical Overlays.
\newblock 2026.
\newblock Preprint in the accompanying series (see prior\_series\_tex/).

\bibitem{bep5}
{BitTorrent}.
\newblock BitTorrent Enhancement Proposal 5: {DHT} Protocol.
\newblock 2008.
\newblock Mainline BitTorrent DHT specification (BEP~5); defines $K=8$ closest nodes returned in responses.
\newblock \url{https://www.bittorrent.org/beps/bep_0005.html}

\bibitem{wang2013mldht}
Liang Wang and Jussi Kangasharju.
\newblock Measuring Large-Scale Distributed Systems: Case of BitTorrent Mainline {DHT}.
\newblock IEEE International Conference on Peer-to-Peer Computing (P2P), 2013.
\newblock DOI: 10.1109/P2P.2013.6688697.
\newblock \url{https://researchportal.helsinki.fi/en/publications/measuring-large-scale-distributed-systems-case-of-bittorrent-main/}

\bibitem{mldhtProjectPage}
Liang Wang and Jussi Kangasharju.
\newblock BitTorrent Mainline {DHT} Measurement (MLDHT) project page.
\newblock 2013.
\newblock University of Helsinki project page; summarizes measurement results (e.g., 15--27 million nodes per day with clear daily churn).
\newblock \url{https://www.cs.helsinki.fi/u/jakangas/MLDHT/}

\bibitem{magnetdb2025}
Scott Seidenberger and Noah Pursell and Anindya Maiti.
\newblock {MagnetDB}: A Longitudinal Torrent Discovery Dataset with IMDb-Matched Movies and TV Shows.
\newblock 2025.
\newblock DOI: 10.48550/arXiv.2501.09275.
\newblock arXiv:2501.09275 (dataset collected via the BitTorrent DHT, 2018--2024)
\newblock \url{https://arxiv.org/abs/2501.09275}

\bibitem{daniel2024ipfs_privacy}
Erik Daniel and Florian Tschorsch.
\newblock Exploring the design space of privacy-enhanced content discovery for Bitswap.
\newblock Computer Communications, 217, 12--24, 2024.
\newblock DOI: 10.1016/j.comcom.2024.01.029.

\bibitem{daniel2023bitswap_privacy}
Erik Daniel and Florian Tschorsch.
\newblock Privacy-Enhanced Content Discovery for {Bitswap}.
\newblock IFIP Networking Conference (IFIP Networking), 1--9, 2023.
\newblock DOI: 10.23919/IFIPNetworking57963.2023.10186387.
\newblock Open-access proceedings PDF via IFIP OpenDL.
\newblock \url{https://opendl.ifip-tc6.org/db/conf/networking/networking2023/1570881323.pdf}

\bibitem{cholez2024ipfs_sybil}
Thibault Cholez and Claudia-Lavinia Ignat.
\newblock Sybil Attack Strikes Again: Denying Content Access in {IPFS} with a Single Computer.
\newblock Proceedings of the 19th International Conference on Availability, Reliability and Security (ARES), 2024.
\newblock DOI: 10.1145/3664476.3664482.
\newblock See also DBLP: conf/IEEEares/CholezI24.

\bibitem{sridhar2024ipfs_censorship}
Srivatsan Sridhar and Onur Ascigil and Navin Keizer and Fran{\c{c}}ois Genon and S{\'e}bastien Pierre and Yiannis Psaras and Etienne Rivi{\`e}re and Micha{\l} Kr{\'o}l.
\newblock Content Censorship in the InterPlanetary File System.
\newblock Network and Distributed System Security Symposium (NDSS), 2024.
\newblock DOI: 10.14722/ndss.2024.23153.
\newblock NDSS 2024; Sybil-based censorship of IPFS provider records and practical mitigations.
\newblock \url{https://www.ndss-symposium.org/wp-content/uploads/2024-153-paper.pdf}

\bibitem{netto2025ipfs_active_sybil}
V. H. M. Netto and Thibault Cholez and Claudia-Lavinia Ignat.
\newblock Active Sybil Attack and Efficient Defense Strategy in {IPFS} {DHT}.
\newblock 2025.
\newblock DOI: 10.48550/arXiv.2505.01139.
\newblock arXiv:2505.01139
\newblock \url{https://arxiv.org/abs/2505.01139}

\bibitem{wang2008attackingkad}
Peng Wang and John A. Tygar and Bernard S. Yee.
\newblock Attacking the {Kad} Network.
\newblock Proceedings of the 1st ACM Workshop on Secure Network Protocols (NPSec), 7--12, 2008.
\newblock DOI: 10.1145/1460877.1460907.

\bibitem{serfling1974sampling}
Serfling, R. J..
\newblock Probability inequalities for the sum in sampling without replacement.
\newblock The Annals of Statistics, 2(1), 39--48, 1974.
\newblock DOI: 10.1214/aos/1176342611.

\bibitem{bardenet2015withoutreplacement}
Bardenet, R{\'e}mi and Maillard, Odalric-Ambrym.
\newblock Concentration inequalities for sampling without replacement.
\newblock Bernoulli, 21(3), 1361--1385, 2015.
\newblock DOI: 10.3150/14-BEJ605.
\newblock Preprint: arXiv:1309.4029.

\bibitem{mazmudar2024peer2pir}
Miti Mazmudar and Shannon Veitch and Rasoul Akhavan Mahdavi.
\newblock Peer2PIR: Private Queries for {IPFS}.
\newblock 2024.
\newblock DOI: 10.48550/arXiv.2405.17307.
\newblock arXiv:2405.17307
\newblock \url{https://arxiv.org/abs/2405.17307}

\bibitem{wang2012sybil}
Liang Wang and Jussi Kangasharju.
\newblock Real-world Sybil Attacks in {BitTorrent} Mainline {DHT}.
\newblock 2012 {IEEE} Global Communications Conference (GLOBECOM), 826--832, 2012.
\newblock DOI: 10.1109/GLOCOM.2012.6503215.
\newblock Unpaywalled author PDF; documents large-scale Sybil presence in Mainline DHT.
\newblock \url{https://nymity.ch/sybilhunting/pdf/Wang2012a.pdf}

\bibitem{ipfs_kad_dht_spec}
{IPFS Project}.
\newblock {IPFS} {K}ademlia {DHT} Specification.
\newblock 2025.
\newblock Version dated 20 November 2025; specifies bucket size/replication parameter $k=20$; recommends client lookup concurrency 10 (denoted $\alpha$ in the spec); and sets provider-record validity to 48h with a recommended republish interval of 22h (plus provider-address TTL 24h).
\newblock \url{https://specs.ipfs.tech/routing/kad-dht/}

\bibitem{i2p_netdb_overview}
{I2P Project}.
\newblock Network Database (netDB): Floodfill storage and lookup.
\newblock 2026.
\newblock Operational documentation: entries are flooded to the three closest floodfills and lookups are typically issued in parallel to two "good" closest floodfills; LeaseSet lookups are garlic encrypted while RouterInfo lookups are not encrypted.
\newblock \url{https://beta.i2p.net/en/docs/overview/network-database/}

\bibitem{karapapas2025ipfs_anonymity_abuse}
Christos Karapapas and Iakovos Pittaras and George C. Polyzos and Constantinos Patsakis.
\newblock Hello, won't you tell me your name?: Investigating Anonymity Abuse in {IPFS}.
\newblock 2025.
\newblock DOI: 10.48550/arXiv.2506.04307.
\newblock arXiv:2506.04307; also published as a chapter in the ARES 2025 International Workshops proceedings volume ``Availability, Reliability and Security'' (Springer), DOI:10.1007/978-3-032-00635-6\_11.
\newblock \url{https://arxiv.org/abs/2506.04307}

\bibitem{mattertran2025ipfs_censorship}
Jan Matter and Muoi Tran.
\newblock Network-level Censorship Attacks in the {I}nter{P}lanetary {F}ile {S}ystem.
\newblock 2025.
\newblock DOI: 10.48550/arXiv.2509.06626.
\newblock arXiv:2509.06626
\newblock \url{https://arxiv.org/abs/2509.06626}


\bibitem{series:oddsinflation}
Anonymous.
\newblock Odds-Inflation Anonymity: A Minimal Endpoint Semantics for Identification Budgets.
\newblock Anonymity series Paper~A, 2026. (This archive.)

\bibitem{series:evalnotary}
Anonymous.
\newblock Notarized Anonymity Certificates: Anti-Grinding Receipts for Evaluation Attempts.
\newblock Evaluation series Paper~1, 2026. (This archive.)

\bibitem{series:release}
Anonymous.
\newblock Release Property, Ledger Receipts, and Public Lineage for Anonymity Claims.
\newblock Release-and-Destination series Paper~A, 2026. (This archive.)

\bibitem{series:artifactpolicy}
Anonymous.
\newblock Artifact policy and supporting materials for the anonymity/AnonDHT series.
\newblock 2026.
\newblock See synthesis paper: Artifact policy (this archive). Artifacts, tooling, certificates, and scripts are available on request as a single archive.
\bibitem{series:contractobjects}
Anonymous.
\newblock Contract Objects for Anonymous-DHT Privacy Budgets and Claims.
\newblock Series synthesis note (Synthesis~0), 2026. (This archive.)

\bibitem{series:abomoinl}
Anonymous.
\newblock ABOM/OINL: Audit BOM and Outcome Identification Noise Label for Anonymous DHT Deployments.
\newblock Anonymity series Paper~C, 2026. (This archive.)

\bibitem{series:traceifaces}
Anonymous.
\newblock Trace Interfaces for Anonymous {DHT}s: Observation Tiers, Committee-Contact Traces, and Exposure Normal Forms.
\newblock Series synthesis note (Synthesis~3), 2026. (This archive.)

\bibitem{series:threatwindows}
Anonymous.
\newblock Threat Windows for Anonymous {DHT}s: What a Budget Covers and For How Long.
\newblock Series synthesis note (Synthesis~4), 2026. (This archive.)


\bibitem{series:notation}
Anonymous.
\newblock Notation and Minimal Model for the One-True-Archive.
\newblock Series synthesis note (Synthesis~8), 2026. (This archive.)

\end{thebibliography}

\end{document}